% ============================================================
% CAP 5 — SÍNTESIS: LOS CINCO INSIGHTS DE DISEÑO
% ============================================================

\section{S\'intesis: los cinco insights de dise\~no}
\label{sec:insights}

Los insights que se presentan a continuación son el puente entre
la investigación de usuario (cap.~\ref{cap:investigacion}) y los
criterios de diseño del servicio (§\ref{sec:qfd}). Cada insight
se estructura en cuatro elementos: la tensión usuario--sistema que lo origina, el dato cuantitativo o cualitativo que lo respalda, la observación que interpreta el patrón, y el criterio de diseño que genera. Ningún criterio de diseño
posterior a este capítulo puede formularse sin que uno de estos
insights lo respalde.

%% ── INSIGHT 1 ──────────────────────────────────────────────

\subsection*{Insight 1 - Fatiga al volante como norma}
\label{sec:insight-1}

\textbf{Tensión:} El sanitario necesita llegar y volver a casa
de forma segura tras 10--12\,h de turno nocturno. El sistema
le obliga a conducir en estado de fatiga porque no existe
alternativa operativa.

\textbf{Dato:} 88 de 167 encuestados (52{,}7\,\%) declaran
conducir con sueño o fatiga \emph{varias veces a la semana}.
El índice de gravedad del ítem es 60{,}2 sobre 100
\parencite{LouridoRegueira2026}. La entrevista en profundidad
confirma este patrón con evidencia cualitativa directa: la
participante declara no esperar el autobús si el tiempo de
espera supera 15~minutos porque «sales reventado de trabajar»,
y refiere conocer compañeras que han sufrido accidentes por
quedarse dormidas al volante de vuelta a casa
(Participante~1, entrevista semiestructurada, mayo de 2026).

\textbf{Observación:} La fatiga al volante post-turno no es un
evento excepcional sino un patrón habitual que el 52{,}7\,\%
de los encuestados experimenta varias veces a la semana. El
sistema de movilidad no ha articulado ninguna alternativa para
este perfil de desplazamiento. La entrevista añade una dimensión
operativa que la encuesta no capturó: la fatiga post-turno
impone un umbral de tolerancia de espera concreto y explícito.
Superado ese umbral, el trabajador descarta el transporte
colectivo y recurre al vehículo privado o al taxi, anulando el
efecto del servicio antes de que este haya tenido oportunidad
de operar.

\textbf{Criterio de diseño derivado:} El servicio debe operar
como transporte de retorno prioritario post-turno. La ventana
crítica de riesgo es la vuelta a casa al amanecer. El tiempo
de espera desde el fichaje de salida no debe superar
15~minutos: superado ese umbral, el trabajador en estado de
fatiga post-turno opta sistemáticamente por alternativas
individuales (Participante~1, entrevista semiestructurada,
mayo de 2026). Este parámetro procede de fuente cualitativa
(n\,=\,1) y requiere validación con muestra representativa
para establecerse como umbral técnico definitivo.

%% ── INSIGHT 2 ──────────────────────────────────────────────

\subsection*{Insight 2 - Aparcamiento: coste oculto sistémico}
\label{sec:insight-2}

\textbf{Tensión:} El trabajador necesita dejar el vehículo cerca
del hospital sin coste ni demora. El sistema le ofrece parking
de personal insuficiente o una opción privada de aproximadamente
80\,€/mes sin plaza garantizada.
\textbf{Dato:} 102 de 167 encuestados (61{,}1\,\%) tienen
dificultades para aparcar \emph{varias veces a la semana}.
45 de 167 (26{,}9\,\%) identifican el aparcamiento garantizado
como el cambio prioritario único \parencite{LouridoRegueira2026}.

\textbf{Observación:} El aparcamiento no es un problema periférico sino la segunda barrera más citada como cambio prioritario. Su coste económico y temporal es invisible en las estadísticas de movilidad pero presente en la experiencia cotidiana del trabajador.

\textbf{Criterio de diseño derivado:} El servicio debe incluir
nodos de embarque en zonas de aparcamiento libre o seguro,
eliminando la necesidad de aparcar en el hospital.

%% ── INSIGHT 3 ──────────────────────────────────────────────

\subsection*{Insight 3 - Transporte público invisible en la franja nocturna}
\label{sec:insight-3}

\textbf{Tensión:} El sanitario necesita una alternativa al coche
puntual y sincronizada con sus turnos. El transporte público no
existe, es desconocido o llega a deshora para el 52{,}1\,\%
de la muestra.

\textbf{Dato:} 87 de 167 encuestados (52{,}1\,\%) carecen de
cobertura fiable de transporte público en la franja nocturna
(39 sin servicio, 48 sin información). 79 de 167 (47{,}3\,\%)
no encuentran transporte público \emph{varias veces a la semana}.
51 de 167 (30{,}5\,\%) señalan la falta de transporte público
como el cambio prioritario \parencite{LouridoRegueira2026}. La
entrevista documenta tres dimensiones adicionales que la encuesta
no capturó: la variabilidad de la hora real de salida por planta
(el turno teórico acaba a las 03:00, pero la participante refiere
salir con frecuencia a las 03:30 por carga de trabajo), la
inoperatividad del servicio en festivos y fines de semana, y la
lentitud de las rutas de autobús disponibles frente al taxi
(aproximadamente 40~minutos con múltiples paradas frente a
10~minutos de trayecto directo) (Participante~1, entrevista
semiestructurada, mayo de 2026).

\textbf{Observación:} El 52{,}1\,\% sin cobertura fiable incluye
tanto la ausencia objetiva de servicio como el desconocimiento
de su existencia. Ambas situaciones producen el mismo efecto
práctico: el trabajador no puede contar con el transporte público
para planificar su desplazamiento nocturno. La entrevista añade
una precisión operativa relevante: la hora de salida real varía
por planta dentro del mismo cambio de turno. Un servicio con una
única salida por cambio de turno excluye a los trabajadores que
terminan con retraso por sobrecarga de planta. La respuesta de
diseño no es solo sincronización con el horario nominal del turno:
es sincronización con la variabilidad real de la salida.

\textbf{Criterio de diseño derivado:} El servicio debe tener
horarios fijos vinculados a los cambios de turno (22:00, 00:00,
07:30), no frecuencias genéricas de red urbana. La arquitectura
de horarios debe contemplar ventanas escalonadas por turno para
absorber la variabilidad real de salida por planta: una única
salida por cambio de turno es insuficiente para dar cobertura a
los trabajadores que terminan con retraso (Participante~1,
entrevista semiestructurada, mayo de 2026). Este parámetro
procede de fuente cualitativa (n\,=\,1) y requiere contraste con
los datos reales de salida por planta del \acs{CHUF} para
dimensionarse como requisito operativo definitivo.

%% ── INSIGHT 4 ──────────────────────────────────────────────

\subsection*{Insight 4 - Demanda espontánea de shuttle}
\label{sec:insight-4}

\textbf{Tensión:} El sanitario necesita un transporte
punto a punto gestionado o coordinado por el sistema sanitario.
Los propios usuarios lo formulan espontáneamente sin que el
cuestionario lo sugiera.

\textbf{Dato:} 21 de 46 respuestas abiertas (45{,}7\,\%)
combinan dos o más problemas simultáneos, señal de necesidad
compuesta. Una respuesta, codificada como T8, propone
espontáneamente un shuttle de empresa sin que el cuestionario
lo induzca. 65 de 167 encuestados (38{,}9\,\%) residen en
un municipio distinto al del hospital
\parencite{LouridoRegueira2026}. La entrevista extiende esta
demanda al perfil conductor: la participante afirma que «mucha
gente que va en coche sí lo usaría», incluso disponiendo de
vehículo propio, e indica que ella misma prescindiría de llevarlo
si existiera el servicio (Participante~1, entrevista
semiestructurada, mayo de 2026).

\textbf{Observación:} La aparición espontánea de la propuesta
de shuttle, sin inducción del cuestionario, indica que la
solución no es contraintuitiva para el usuario: existe una
demanda latente que el sistema sanitario no ha articulado pero
que los propios trabajadores ya han formulado. La entrevista
precisa tres requisitos operativos que los usuarios asocian a
un servicio creíble: horarios adaptados a la hora real de salida
y no al horario nominal del turno, acceso exclusivo para
personal sanitario acreditado como condición de seguridad
percibida, y gestión con participación del hospital porque
«saben cómo funciona un hospital» y conocen la variabilidad
operativa real (Participante~1, entrevista semiestructurada,
mayo de 2026).

\begin{quote}
«Estaría bien que a los trabajadores de hospitales se nos
pusiera un transporte para llevarnos al hospital y recogernos
en un punto y de ahí para casa.»
\hfill\textit{Respuesta abierta, código T8 \parencite{LouridoRegueira2026}}
\end{quote}

\textbf{Criterio de diseño derivado:} El servicio debe
articularse como transporte punto a punto entre nodos
periféricos y el hospital, no como extensión de línea urbana.
La arquitectura operativa debe incorporar tres condiciones
derivadas de la investigación cualitativa: acceso escalonado
que absorba la variabilidad real de salida por planta,
exclusividad de uso para personal sanitario acreditado, y
participación del hospital o del \acs{SERGAS} en la gobernanza
del servicio (Participante~1, entrevista semiestructurada,
mayo de 2026). Estos tres parámetros proceden de fuente
cualitativa (n\,=\,1) y requieren validación con muestra
representativa antes de establecerse como requisitos operativos
definitivos.

%% ── INSIGHT 5 ──────────────────────────────────────────────

\subsection*{Insight 5 - Brecha de género en el riesgo}
\label{sec:insight-5}

\textbf{Tensión:} Las mujeres sanitarias absorben de forma
desproporcionada el riesgo de fatiga al volante. El sistema
impone el mismo riesgo a todos, pero la carga recae
mayoritariamente sobre las mujeres: el 90{,}8\,\% de ellas
conduce fatigada con una frecuencia de al menos varias veces
al mes, frente al 53{,}8\,\% de los hombres.

\textbf{Dato:} 138 de 152 mujeres de la muestra (90{,}8\,\%)
presentan riesgo de fatiga al volante con frecuencia igual o
superior a \emph{varias veces al mes}, frente a 7 de 13 hombres
(53{,}8\,\%). La muestra es 91{,}0\,\% femenina (152 de 167)
\parencite{LouridoRegueira2026}. La entrevista en profundidad
aporta una jerarquía de prioridades formulada explícitamente
por la participante al evaluar un servicio de transporte
nocturno: seguridad, rapidez, independencia y coste, en ese
orden (Participante~1, entrevista semiestructurada, mayo de 2026).
La seguridad es el primer criterio de elección modal, por encima
de la velocidad o el coste.

\textbf{Observación:} La brecha de 37 puntos porcentuales entre
mujeres (90{,}8\,\%) y hombres (53{,}8\,\%) en riesgo de fatiga
al volante, en una muestra 91\,\% femenina, significa que el
riesgo de conducción nocturna bajo fatiga es estadísticamente
un problema de género, aunque el sistema lo trate como un riesgo
individual. La entrevista precisa la localización espacial de
ese riesgo: la inseguridad percibida no se produce durante el
trayecto a bordo, sino en la espera en parada de vía pública.
La participante declara llamar a su madre por teléfono mientras
espera en la parada y refiere que el malestar desaparece en
cuanto sube al vehículo. El punto de diseño crítico no es el
interior del vehículo: es el tiempo y el espacio de espera
previos al embarque (Participante~1, entrevista semiestructurada,
mayo de 2026).

\textbf{Criterio de diseño derivado:} El servicio debe garantizar
accesibilidad y seguridad percibida para usuarias que viajan
solas en franja nocturna. El diseño del espacio de espera tiene
prioridad sobre el diseño del interior del vehículo: el riesgo
percibido se concentra en la parada, no en el trayecto. Los
nodos de embarque deben ubicarse en recintos cubiertos y
protegidos, no en paradas de vía pública convencional
(Participante~1, entrevista semiestructurada, mayo de 2026).
Este parámetro procede de fuente cualitativa (n\,=\,1) y
requiere validación con muestra representativa para establecerse
como requisito de localización definitivo.
